% Gerado de docs/monografia/08-conclusao.md por scripts/monografia_latex.py.
% Não edite aqui: edite o .md e gere de novo.

\chapter{\texorpdfstring{\mbox{CONCLUSÃO}}{CONCLUSÃO}}\label{cap:8}

Este trabalho partiu de uma dificuldade conhecida do ensino de ponteiros em C: uma resposta incorreta não distingue o estudante que sustenta uma concepção errônea daquele que desconhece o conteúdo, e a concepção errônea não é corrigida apenas com explicações adicionais. Para tratá-la, foi desenvolvido o Structura, sistema que combina dois mecanismos. O diagnóstico utiliza a resposta, a razão e a certeza declarada para formular, a cada premissa, uma hipótese sobre o estudante: se sabe, se não sabe ou se sustenta uma concepção errônea. A intervenção decorre do diagnóstico: ao estudante que não sabe, é oferecida uma escada de passos curtos; ao que sustenta uma concepção errônea, uma demonstração na qual ele prevê o comportamento de um programa e observa o resultado da execução.

O objetivo geral, de desenvolver e verificar um sistema com essas características, foi alcançado, e os objetivos específicos foram atendidos da seguinte forma. O conteúdo foi organizado em premissas encadeadas, com cada uma das 61 concepções errôneas do catálogo associada à premissa que nega e com procedência declarada (Seções~3 e 4). O diagnóstico e a escolha do percurso foram implementados conforme descrito na Seção~3. Os 13 tópicos, com 79 premissas, 402 passos e 80 demonstrações, além das 79 cenas que animam a revisão de cada premissa, foram conferidos contra o compilador: o que cada passo, cada demonstração e cada cena afirma sobre um programa corresponde ao comportamento do programa no ambiente especificado (Seções~5 e 6). No banco de questões, a conferência por compilador alcança as 69 questões de prever a saída com resposta em valor, e as demais dependem da revisão automatizada da forma e de leitura (Seções~5.2 e 7.1). O comportamento foi conferido com dez perfis de estudante simulado em todos os tópicos: cada caso segue o percurso especificado, e a extensão da sessão é limitada.

A questão de pesquisa indagava se é possível construir esse sistema de modo que o conteúdo e o comportamento possam ser conferidos sem depender do julgamento de quem os elaborou. Nos limites do que foi realizado, a resposta é afirmativa quanto ao que o compilador e os testes alcançam: o comportamento dos programas no ambiente especificado e o comportamento do sistema. A verificação automatizada mostrou-se útil ao longo do processo, ao identificar defeitos que a leitura não havia percebido, entre eles uma questão cuja resposta indicada como correta não é válida no compilador utilizado na verificação e oito questões cujo código era exibido de forma truncada. A leitura do conteúdo, iniciada pelo autor e estendida aos demais passos com assistência do modelo de linguagem, identificou, por sua vez, o que a verificação automatizada não poderia identificar: perguntas verdadeiras, mas que repetiam a frase anterior, reescritas em 101 dos 244 passos guiados. As duas formas de verificação são complementares, e nenhuma substitui a outra.

O trabalho não responde se o sistema produz aprendizagem, nem se a hipótese diagnóstica corresponde ao que o estudante de fato pensa. Essas questões requerem um estudo com estudantes, que constitui o primeiro dos trabalhos futuros: o protocolo está definido (Apêndice~A), e o sistema já registra os dados necessários e os exporta em tabelas, pelo código do participante (Seção~3.7 e Apêndice~C). A adequação pedagógica do conteúdo didático, elaborado com assistência de um modelo de linguagem e validado pelo autor, também depende desse estudo. As contribuições deste trabalho são, portanto, as relacionadas na Seção~1.7: um catálogo de concepções errôneas com procedência declarada, um diagnóstico que utiliza a certeza declarada para escolher o percurso, um conteúdo didático conferido por compilador e um método no qual as afirmações do material de ensino sobre o comportamento dos programas são verificadas por compilador, e não pelo julgamento de quem as redigiu.
